\section{最大值原理：把最优轨线变成两点边值}
\label{sec:16-Control-and-Reinforcement-Learning/01-Optimal-Control/02}

上一节解决了``稳住''——偏差来了，有办法把它拉回来。但稳住只是及格线。到目标的路径有无数条：有的费时间，有的费燃料，有的把人颠散。如何在无数条路径里挑出代价最小的那一条？这一节把``最优''二字形式化，然后用一个思想的连续转身，把函数空间上的寻优问题压成一组微分方程加边界条件。

\subsection{先有数字：一辆单车、一列火车、二十分钟}

你骑一辆单车赶一班还有 20 分钟发车的列车。车站距离你 5 公里，当前车速 15 km/h。任务很简单：选一条蹬车策略 \(u(t)\)——踩多重、多快、在哪一段歇——在 20 分钟之内到达站台。

这件事有四个零件需要写清楚。

\textbf{状态方程。}你的位置 \(x_1\) 和速度 \(x_2\) 怎样响应蹬车力度 \(u\)。可以简单假定 \(\dot x_1 = x_2\)，\(\dot x_2 = u - \gamma x_2\)（阻力与速度成正比），初值 \(x(0) = x_0\) 已知。这里的动力学不一定只是牛顿定律——任何物理过程只要写成 \(\dot x = f(x, u)\)，都符合后面所有推导的接口。

\textbf{控制的可行范围。}腿不会无限输出：\(u(t) \in [u_{\min}, u_{\max}]\)，也许弯道处还要额外限制。这是你被物理世界绑住的手。

\textbf{终端代价。}发车时没到站——迟到代价很大。冲得太猛来不及停——冲出站台的代价也很大。终端代价 \(\Phi(x(T))\) 把``终点状态的好坏''兑现成数字。

\textbf{运行代价。}路上蹬累了、加太急减速不舒服——\(L(x, u)\) 把沿途的消耗和不适积起来。

目标函数：

\[
J(u) = \Phi(x(T)) + \int_0^T L(x(t), u(t))\, dt.
\]

在所有满足约束的 \(u(\cdot)\) 中找那个让 \(J\) 最小的。

注意一件事：\(u(\cdot)\) 是一整条函数，不是一个数。这是个无穷维空间上的优化——你微积分课里学的``求导令零解方程''不能直接搬过来，因为自变量不再是 \(x \in \R\) 而是一个函数。需要新的梯度和新的最优性条件。

\subsection{第一次尝试：把时间剁碎}

一个很自然的想法——既然计算机只能处理离散，那把连续时间切成 \(N\) 小段，在离散格子里用动态规划，让 \(N\) 越切越细趋近连续解。这是 Bellman 那条线的思路（下一节细讲）。

但它马上撞到规模墙。每一步都要在状态空间的所有点求值函数——状态空间稍微高维，离散点的数目就爆炸。单车的位置和速度只是两维，还能应付；换成一架四轴飞行器的 12 维状态空间，网格的格点数随维数按指数增长。这条路是对的，但代价很重。

能不能不遍历整个状态空间，只在一整条候选轨线上找到最优性应该满足的条件？

\subsection{换成变分的思路：扰动最优轨线}

设 \(u^*\) 是一条最优控制轨线，\(x^*\) 是对应的状态轨线。最优的意思是：任何微小扰动都不能让代价再降。

对 \(u^*\) 施加一个微扰 \(u^* + \varepsilon \eta(t)\)，状态相应地偏移 \(x^* + \varepsilon \delta x(t) + o(\varepsilon)\)。代入动力学线性化：

\[
\dot{\delta x} = f_x \delta x + f_u \eta, \quad \delta x(0) = 0,
\]

其中 \(f_x\) 和 \(f_u\) 都是沿最优轨线取值的 Jacobian。

代价的一阶变分为

\[
\delta J = \frac{\partial \Phi}{\partial x}(x^*(T)) \cdot \delta x(T) + \int_0^T \left( L_x \cdot \delta x + L_u \cdot \eta \right) dt.
\]

这里有件棘手的事：\(\eta\) 是你亲手加的扰动，你完全掌控它的形状。但 \(\delta x\) 夹在积分里——它不是你直接控制的量，它需要积微分方程才能算出来。你好像在跟一个看不到的影子打架：想调 \(\eta\) 让 \(\delta J\) 降，但 \(\delta x\) 的影响绕了远路，直接条件写不干净。

\subsection{一个看似多此一举的操作：乘上任意函数再积分}

沿任何轨线，方程 \(\dot x - f(x,u) = 0\) 在任何时刻都成立。随便取一个函数 \(\lambda(t)\)，乘上去、从 0 积到 \(T\)：

\[
\int_0^T \lambda(t) \cdot (\dot x - f(x,u))\, dt = 0.
\]

加进 \(J\) 不改变\(J\) 的数值——你的全是零。但这个恒等式给了你来回搬导数的自由度。

具体地，把这个零加进去后，用分部积分把 \(\lambda \cdot \delta \dot x\) 拆成边界项减 \(\dot\lambda \cdot \delta x\)。加上原来的 \(\delta J\) 中的各项，\(\delta x\) 前面的系数汇集为几个表达式。这时候你做一件事：你自由选择 \(\lambda(t)\) 的演化方程和终端条件，专门让 \(\delta x\) 的所有系数恰好消为零。

一旦 \(\delta x\) 被消掉，剩下的被积项只含 \(\eta\)——而 \(\eta\) 是任意函数，它前面的系数必须恒为零才能保证最优性。这个条件就给出了最优控制的逐点特征。

这就是协态（costate，也译伴随变量、adjoint）的来源。它不是物理量——它是你为了把 \(\delta x\) 从必要条件里解耦掉而引入的辅助函数。数学家有时叫它 adjoint，是``伴随方程''的伴随。协态的引入看起来像数学技巧，但它的效果截然不同：把``\(\delta x\) 需要解方程才知道''这个障碍彻底铲掉。

\subsection{把结论写干净}

沿着上述变分推导，三块拼成一套完整的条件。定义 Hamiltonian：

\[
H(x, u, \lambda) = \lambda \cdot f(x, u) - L(x, u).
\]

\textbf{Pontryagin 最大值原理}（写出必要条件的形式）。若 \(u^*\) 是最优控制，\(x^*\) 为对应的状态轨线，则存在协态函数 \(\lambda(\cdot)\)，满足

伴随方程：
\[
\dot\lambda = -\frac{\partial H}{\partial x},
\]

终端条件：
\[
\lambda(T) = -\frac{\partial \Phi}{\partial x}(x^*(T)),
\]

并且在每一时刻 \(t\)：
\[
u^*(t) \in \argmax_{u \in U} H(x^*(t), u, \lambda(t)).
\]

符号有一处要交代。另一种盛行的约定是 \(H = L + \lambda \cdot f\)，然后每时刻取最小，协态差个符号。两种陈述数学上完全等价。本节依照``最大值''的名字沿用取最大的约定；下一节切换到后一种约定后，协态恰好成为值函数的梯度，读法上更自然。

这个定理的重心在结构转换。``整条时间轴上的全局寻优''——无穷维——被压成了``每一时刻在 \(U\) 上的逐点寻优''——有限维，\(U\) 常常是 \([u_{\min}, u_{\max}]\) 之类的区间，有时甚至可以直接读出答案。代价是必须同时解出 \(x^*\) 和 \(\lambda\)，而这两条轨线互相缠绕。

\subsection{协态怎么读：状态的边际价值}

上面说协态是为``消去 \(\delta x\)''而设的技巧。如果只看它的定义，会误以为它只是一个没有物理含义的中间变量。换个记号约定之后能看清它的经济含义。

改取 \(H = L + \lambda \cdot f\) 和逐点最小化的约定（协态变号）。终端条件成为 \(\lambda(T) = \partial \Phi / \partial x\)——没有负号。在这个约定下：

\(\lambda(t)\) 是这样一个量：你在时刻 \(t\) 给状态一个微小的正推动，剩下的最优总代价会改变多少。它是状态对剩余目标的边际价值。这和约束优化中拉格朗日乘子被读作影子价格完全一致——约束多松一单位，目标改善多少。

伴随方程 \(\dot\lambda = -\partial H / \partial x\) 就说：这个边际价值沿最优轨线按``当前状态对远期代价的边际影响''的速率折旧。如果状态对远期代价的正影响大，边际价值折旧就快；否则折旧就慢。

最大化 Hamiltonian 的第一项 \(\lambda \cdot f\) 可以这么读：选控制 \(u\) 时，不只看眼下的运行代价 \(L\)——还要看这个控制把状态推向哪个方向。在边际价值的度量下，往值钱的方向推就加分，往赔钱的方向推就扣分。

\subsection{落成两点边值问题}

拼起来看骨架：

\begin{itemize}
\tightlist
\item 状态方程 \(\dot x = f\)——从左端初值 \(x(0) = x_0\) 向前积分；
\item 伴随方程 \(\dot\lambda = -\partial H / \partial x\)——从右端终端条件 \(\lambda(T) = -\partial \Phi / \partial x\) 向后积分；
\item 每时刻的 \(u\) 由逐点最大化同时约束两者。
\end{itemize}

条件分居时间轴两端：\(x(0)\) 已知，\(\lambda(T)\) 由终端代价决定，但 \(\lambda(0)\) 未知，\(x(T)\) 也未知。没法从头一口气积到底。

这恰是 \hyperref[sec:08-Numerical-Analysis/06-ODE/04]{``边值问题：从打靶到全局离散''} 的研究对象。打靶法猜一个 \(\lambda(0)\)，正向积分至 \(T\)，逐点最大化给出 \(u(t)\)，看终端偏差，修正猜测再试。多重打靶把时间轴切成多段，每段以自己的初值做变量；配置法则把轨线的整个形态离散为多项式后一次性求解。

最优控制把``寻优''的困难转化成了``解方程''的困难，再转交给数值方法。每一转都有代价，但每一转都换来了可计算性。

\subsection{Bang-bang：不解方程先读出解的形状}

最大值原理有时给的是解的\textbf{形状}，而不是具体的数值轨线。经典案例：最速停车。

动力学 \(\ddot y = u\)，控制受限于 \(\abs{u} \le 1\)。要求从初始位置以最短时间到达原点并停住。运行代价只计时间：\(L \equiv 1\)。Hamiltonian：

\[
H = \lambda_1 \dot y + \lambda_2 u - 1.
\]

\(H\) 对 \(u\) 是线性的。在 \([-1, 1]\) 上求最大——线性函数的极值在端点：

\[
u^*(t) = \sign \lambda_2(t).
\]

控制只取 \(\pm 1\) 两个极端，永不取中间值——要么全油门加速，要么全力刹车。Bang-bang 控制的名字来源于此。

再看伴随方程：\(\dot\lambda_1 = 0\)，\(\dot\lambda_2 = -\lambda_1\)。\(\lambda_2\) 是一次函数，至多跨零一次。切换至多一次：先满推一段，再满刹一段。

从头到尾没有积任何具体的初值问题——最长切换次数、控制的极端性，这些结构信息直接从最大值原理里落出来了。``解的形状先于解本身''——这是逐点最大化带来的独有红利。

\subsection{这条路的几个边界}

最大值原理给的是必要条件。满足条件的轨线叫极值轨线（extremal）——它可能是最优，也可能只是临界轨道。非凸动力学或代价之下可能有多条极值轨线，需要横向比较候选解，或者补充充分性条件才能断言全局最优。

协态为什么从终点倒着走？因为它度量的是``当前状态对剩余未来代价的边际影响''。未来的终端代价先定出边界值，这个边际价值再沿时间反向传播到当前。如果终点没有额外约束且终端代价为零，自然终端条件退化为 \(\lambda(T) = 0\)。如果终点被硬性约束扣在某个精确值上，终端条件还会多出新的乘子，不能直接把自由终点的公式套上去。

从理论上说，还有一个更隐蔽的坑：最大值原理的原始陈述要求 \(u\) 取值于一个紧集，状态方程对 \(u\) 可微。对大多数工程问题这两个条件成立。但一旦涉及到脉冲控制、奇异弧（切换函数在非零区间上恒为零导致逐点最大化暂时不能确定 \(u\)），就需要回到更精密的变分分析。

最大值原理是一套强大的结构分析工具——它把无穷维寻优压成了可计算的微分方程加逐点代数问题。但它不替你验证全局最优性，也不替你处理硬约束的乘子结构和奇异弧。知道这些边界，才能在不同的具体问题中恰当地拉满它的力量、也恰当地止住过分的使用。
